\chapter{The master arm}
\label{ch:master}

The master is a miniature mannequin reproducing the Gen3's joint sequence at
roughly a third of its scale, moved with the fingers of one or both hands.
Source: \repo{firmware/master\_arm/teensy\_final.ino}.

\section{Mechanical chain and electronics}

Each arm carries seven potentiometer-instrumented joints, alternating a roll
about the preceding link's own axis with a bend perpendicular to it; link
lengths sum to \SI{0.272}{\metre}. A roll joint moves the tip only through a
downstream bend (rotating a straight arm about its own axis translates the
tip by zero), which is why the loss of J2 or J4 costs reach while J3 or J5
costs almost nothing.

\begin{figure}[htbp]
  \centering
  \includegraphics[width=0.85\linewidth]{master_arm_cad_views.png}
  \caption[The master arm, as designed]{Four orthographic views of the
  assembly at the zero configuration. The mechanical model carries seven
  potentiometer housings, a shaft coupling at each joint, and a wrist
  platform for the inertial sensor and force-sensitive resistor.}
  \label{fig:master-cad}
\end{figure}

A Teensy 4.1 reads fourteen potentiometers, two MPU-6050 inertial sensors,
two force-sensitive resistors and two buttons, emitting one ASCII text frame
per cycle at \SI{100}{\hertz} (\cref{tab:frame}). An earlier firmware
revision placed two potentiometer channels on \texttt{A4}/\texttt{A5}, which
are physically the default \texttt{Wire} bus both inertial sensors use;
calling \texttt{analogRead()} on a pin the I\textsuperscript{2}C peripheral
is driving corrupts the bus intermittently and returns meaningless
potentiometer values, with no symptom that looks like a pin conflict from
the outside. The fix excludes both pins from analog use: Teensy 4.1 exposes
eighteen analog-capable pins, minus two leaves sixteen, exactly what
fourteen potentiometers plus two force sensors require -- no spare.

\begin{table}[htbp]
  \centering
  \small
  \caption{Serial frame, in transmission order.}
  \label{tab:frame}
  \begin{tabular}{@{}llll@{}}
    \toprule
    Field & Count & Format & Meaning \\
    \midrule
    \texttt{k1j1..k1j7} & 7 & \texttt{\%.1f} & left arm joints, \SIrange{0}{360}{\degree} \\
    \texttt{k2j1..k2j7} & 7 & \texttt{\%.1f} & right arm joints \\
    K1IMU, K2IMU & 6+6 & \texttt{\%d} & $a_x a_y a_z g_x g_y g_z$, raw LSB \\
    fsr1, fsr2 & 2 & \texttt{\%d} & grip force, raw counts \\
    btn1, btn2 & 2 & \texttt{\%d} & latched toggle state \\
    \bottomrule
  \end{tabular}
\end{table}

\section{Signal conditioning}

Each potentiometer channel is read at 12-bit resolution and passed through
an integer exponential moving average,
\begin{equation}
  e_{n+1} = \frac{\alpha r_n + (256-\alpha) e_n}{256}, \quad \alpha=64,
  \label{eq:ema}
\end{equation}
a smoothing factor of 0.25 per sample. The filtered count maps to degrees by
a clamped linear transform that discards the outer 50 counts at each end of
the ADC range, where a potentiometer wiper is least linear.

Recorded captures show each potentiometer producing a new \emph{distinct}
value at only \SIrange{14.7}{17}{\hertz} against the \SI{100}{\hertz} frame
rate, so \SIrange{82}{88}{\percent} of consecutive rows repeat -- every
per-row statistic computed before this was noticed was invalidated by it
(\cref{ch:instrument}). At the power-on default of $\pm\SI{2}{\gram}$ full
scale, accelerometer sensitivity is \SI{16384}{LSB\per\gram}; the widely
copied value of \SI{8192}{LSB\per\gram} corresponds to the $\pm4g$ range and,
used here, halves every reported acceleration -- invisible in the elevation
estimate (a uniform scale error cancels in a normalised vector) but fatal to
the accelerometer gate, which then rejects every frame and looks like a dead
sensor. Force-sensitive resistors are read as raw integers with no
filtering; the mapping applies its own deadband and hysteresis downstream,
so the discrepancy between the design (which specifies a filtered value)
and the firmware in hand is recorded rather than resolved.

\section{The CAD is not the kinematic source}
\label{sec:cad-not-source}

Neither STEP file contains a single kinematic-joint entity; they are solid
geometry for manufacture, describing one fixed configuration. No URDF or
xacro anywhere in the workspace references either file -- \textbf{the master
arm has no URDF at all}. What the software uses is seven measured constants
in \repo{srl\_teleop/master\_calibration.py}, taken directly from the
physical arm. The CAD's joint-module pitch agrees with those constants to
\SI{1.2}{\percent}, which is a cross-check, not a licence to treat the CAD
as the source: it cannot distinguish the \SI{43}{\milli\metre} roll from the
\SI{37}{\milli\metre} bend within one measured pitch.
